\SeventhArticle{19}
Since it suffices to know either of the functions $T,U$, we preferred the determination of the latter as the simpler course. Just as this is based on expanding the series $u^{-1}+\frac13u^{-3}+\frac15u^{-5}+\text{etc.}$ into a continued fraction, by similar reasoning we could have used the expansion of the series $t^{-1}+\frac12t^{-2}+\frac13t^{-3}+\frac14t^{-4}+\text{etc.}$ into the continued fraction
\SourceLeaf{199}
\[
\cfrac1{t-\cfrac{\frac12}{1-\cfrac{\frac16}
{t-\cfrac{\frac26}{1-\cfrac{\frac2{10}}
{t-\cfrac{\frac3{10}}{1-\mathrm{etc.}}}}}}}
\]
to derive an algorithm for determining the function $T$ for successive values of the number $n$. We reach the same conclusion, however, by considering that $T$ is nothing other than $\frac U{2^{n+1}}$, or $\frac{W^{(n+1)}}{2^{n+1}}$, if $u$ is replaced by $2t-1$. In this way the functions to be taken successively for $T$ are obtained by the following algorithm:
\[
\begin{aligned}
W&=1\\
\tfrac12W'&=t-\tfrac12\\
\tfrac14W''&=(t-\tfrac12).\tfrac12W'-\frac{1.1}{2.6}W
\\[3pt]&\qquad=tt-t+\tfrac16\\
\tfrac18W'''&=(t-\tfrac12).\tfrac14W''-\frac{2.2}{6.10}.\tfrac12W'
\\[3pt]&\qquad=t^3-\tfrac32tt+\tfrac35t-\tfrac1{20}\\
\tfrac1{10}W''''&=(t-\tfrac12).\tfrac18W'''-\frac{3.3}{10.14}.\tfrac14W''
\\[3pt]&\qquad=t^4-2t^3+\tfrac97tt-\tfrac27t+\tfrac1{70}
\end{aligned}
\]
etc. By induction, this gives in general
\[
\begin{aligned}
T={}&t^{n+1}-\frac{(n+1)^2}{1.(2n+2)}t^n
+\frac{(n+1)^2.nn}{1.2.(2n+2)(2n+1)}t^{n-1}\\
&-\frac{(n+1)^2.nn.(n-1)^2}{1.2.3.(2n+2)(2n+1).2n}t^{n-2}
+\mathrm{etc.}
\end{aligned}
\]
or $T=t^{n+1}F(-(n+1),-(n+1),-2(n+1),t^{-1})$; this induction can easily be given the force of a proof. If one prefers, $T$ can also be expressed with the order of the terms reversed, by
\[
\pm\frac{1.2.3.4\mathinner{.\,.\,.\,.\,.}(n+1)}
{2.6.10.14\mathinner{.\,.\,.\,.\,.}(4n+2)}
F(n+2,-(n+1),1,t)
\]
where the upper sign holds for odd $n$, the lower for even. Finally, in a similar way, $T''$ is found in general to be equal to the product of
\[
\frac{1.1.2.2.3.3\mathinner{.\,.\,.\,.\,.}(n+1).(n+1)}
{2.6.6.10.10.14\mathinner{.\,.\,.\,.\,.}(4n+2).(4n+6)}t^{-(n+2)}
\]
and the infinite series
\[
\begin{aligned}
1&+\frac{(n+2)^2}{1.(2n+4)}t^{-1}
+\frac{(n+2)^2(n+3)^2}{1.2.(2n+4)(2n+5)}t^{-2}\\
&+\frac{(n+2)^2.(n+3)^2(n+4)^2}
{1.2.3.(2n+4)(2n+5)(2n+6)}t^{-3}+\mathrm{etc.}
\end{aligned}
\]
or $F(n+2,n+2,2n+4,t^{-1})$
\SourceLeaf{200}
\SeventhArticle{20}
Since in the function $U$ the powers $u^n$, $u^{n-2}$, $u^{n-4}$, etc., are absent, the roots of the equation $U=0$ will always occur in pairs of equal magnitude and opposite sign; for an even value of $n$, the single root $0$ must also be added. Once the roots have been found, the values of the coefficients $R,R',R''$, etc., will be given, by the method of Article 11, by a polynomial in $u$, which, for an odd value of $n$, will have the form
\[
\gamma u^{n-1}+\gamma'u^{n-3}+\gamma''u^{n-5}+\mathrm{etc.}
\]
but for an even value, if the coefficient corresponding to the root $u=0$ is excluded, the form
\[
\gamma u^{n-2}+\gamma'u^{n-4}+\gamma''u^{n-6}+\mathrm{etc.}
\]
The example of Article 12 shows this very reduction for $n=6$. Evidently, therefore, opposite values of $u$ always correspond to equal coefficients. Further, in the case where $n$ is even, the coefficient corresponding to the root $u=0$ can easily be assigned in general a priori. This coefficient is obtained by substituting in $\frac{U'}{(\frac{\mathrm{d}U}{\mathrm{d}u})}$ the value $u=0$. We have already given the value of the numerator $U'$ for $u=0$ in Article 18; the value of the denominator, however, will be given by the same article as
\[
\begin{aligned}&=\pm\frac{3.5.7\mathinner{.\,.\,.\,.\,.}(n+1)}
{(n+3)(n+5)\mathinner{.\,.\,.\,.\,.}(2n+1)}\\[5pt]&=\pm\frac{3.3.5.5.7.7\mathinner{.\,.\,.\,.\,.}(n+1)(n+1)}
{3.5.7.9.11\mathinner{.\,.\,.\,.\,.}(2n+1)}\end{aligned}
\]
and hence the required coefficient is
\[
=\left(\frac{2.4.6.8\mathinner{.\,.\,.\,.\,.} n}
{3.5.7.9\mathinner{.\,.\,.\,.\,.}(n+1)}\right)^2
\]

\SeventhArticle{21}
The polynomial in $u$ that represents the coefficients $R,R',R''$, etc., in the case we are considering can also be found independently of the general method of Article 11, as follows. By differentiating the equation
\[
\varphi-\frac{U'}U=\frac{U''}U
\]
then substituting $\frac{\mathrm{d}\varphi}{\mathrm{d}u}
=\frac1{1-uu}$, and multiplying by $UU(uu-1)$, we obtain
\[
\begin{aligned}&(uu-1)U'\frac{\mathrm{d}U}{\mathrm{d}u}
-U\left(\frac{\mathrm{d}U'}{\mathrm{d}u}.(uu-1)+U\right)\\[5pt]&=(uu-1)UU\frac{\mathrm{d}(\frac{U''}U)}{\mathrm{d}u}\end{aligned}
\]
\SourceLeaf{201}
The terms on the left of this equation evidently form a polynomial in $u$; therefore, on the right, the coefficients of powers of $u$ with negative exponents must necessarily cancel one another.

But $\frac{\mathrm{d}\frac{U''}U}{\mathrm{d}u}$ gives an infinite series beginning with the term
\[
-\left(\frac{1.2.3.4\mathinner{.\,.\,.\,.\,.}(n+1)}
{1.3.5.7\mathinner{.\,.\,.\,.\,.}(2n+1)}\right)^2u^{-(2n+4)}
\]
and hence, when this series is multiplied by $(uu-1)UU$, nothing can result except the constant quantity
\[
-\left(\frac{1.2.3.4\mathinner{.\,.\,.\,.\,.}(n+1)}
{1.3.5.7\mathinner{.\,.\,.\,.\,.}(2n+1)}\right)^2
\]
Hence we conclude\HistoricalNote{20} that
\[
(uu-1)U'\frac{\mathrm{d}U}{\mathrm{d}u}
+\left(\frac{1.2.3.4\mathinner{.\,.\,.\,.\,.}(n+1)}
{1.3.5.7\mathinner{.\,.\,.\,.\,.}(2n+1)}\right)^2
\]
is divisible by $U$; therefore, to the rational function $\frac{U'}{(\frac{\mathrm{d}U}{\mathrm{d}u})}$, which supplies the coefficients $R,R',R''$, etc., the following polynomial will be equivalent:
\[
-\left(\frac{1.3.5.7\mathinner{.\,.\,.\,.\,.}(2n+1)}
{1.2.3.4\mathinner{.\,.\,.\,.\,.}(n+1)}U'\right)^2.(uu-1)
\]
Instead of this function, which has degree $2n+2$ and evidently involves only even powers of $u$, the remainder from its division by $U$ may be taken; this will have degree $n$ or $n-1$, according as $n$ is even or odd. If, however, in the former case we prefer to exclude the coefficient corresponding to the root $u=0$, we shall take, instead of that function, its remainder from division by $\frac Uu$, which will rise only to degree $n-2$.

